\section{Constant accuracy and the main lower bound}\label{h:end:section}

Rounding the horizon up to a whole number of blocks loses less than one
block. At fine accuracy this loss is negligible. For the remaining
accuracies, even an estimator retaining all the data needs a linear
number of equations. Reflection across their row span proves this
without identifying the full conditional law of the signal.

\begin{lemma}[Reflection with full data]
\label{h:ev:reflection}
Let \(d\ge1\) be an integer and \(T\ge0\). Let \(S\) be uniform on
\(S^{d-1}\). Let a matrix \(A\) with \(d\) columns, independent of
\(S\), have row span \(U\) of dimension at most \(T\). Suppose an
output \(\widehat S\) is drawn from a probability kernel depending
only on \((A,AS)\). Put \(u=P_US\) and \(v=(I-P_U)S\).
For an arbitrary output in \(\mathbb R^d\) and \(0<\epsilon<1\),
\begin{equation}\label{h:ev:direct-residual}
 \Pr\{\|\widehat S-S\|\le\epsilon\}
 \le\frac12+\frac{T}{2d(1-\epsilon^2)}.
\end{equation}
More generally, for \(0<b_0<1\) and
\(0\le\epsilon<\sqrt{1-b_0}\),
\begin{equation}\label{h:ev:antipodal}
 \Pr\{\|\widehat S-S\|\le\epsilon\}
 \le\frac12+\frac12\Pr\{\|u\|^2\ge b_0\}
 \le\frac12+\frac{T}{2b_0d}.
\end{equation}
For a unit output one also has
\begin{equation}\label{h:ev:reflection-expectation}
 \mathbb E\langle\widehat S,v\rangle=0,\qquad
 \mathbb E\langle\widehat S,S\rangle
 \le\mathbb E\|u\|\le\sqrt{T/d}.
\end{equation}
\end{lemma}

\begin{proof}
The projection \(P_U\) is a Borel function of \(A\), since
\(P_U=\lim_{\eta\downarrow0}A^{\mathsf T}(AA^{\mathsf T}+\eta I)^{-1}A\);
for a matrix with no rows it is zero.
Conditional on \(A\), reflection across \(U\) sends \(S=u+v\) to
\(u-v\). It preserves uniform spherical measure and the data \((A,AS)\),
because \(Av=0\), so the output kernel is the same at both signals.
If \(\|v\|>\epsilon\), the reflected signals have distance greater
than \(2\epsilon\); a fixed estimate can succeed for at most one of
them. Averaging the common output kernel over the reflected pair gives
success at most one half on this event. On its complement use the
bound one. Therefore
\[
 \Pr\{\|\widehat S-S\|\le\epsilon\}
 \le\frac12+\frac12\Pr\{\|u\|^2\ge1-\epsilon^2\}.
\]
Spherical isotropy and independence of \(A\) give
\[
 \mathbb E(\|u\|^2\mid A)=\frac{\dim U}{d}\le\frac Td.
\]
Markov's inequality proves \eqref{h:ev:direct-residual}. The same
pairing on the event \(\|u\|^2<b_0\), where
\(\|v\|>\sqrt{1-b_0}>\epsilon\), proves
\eqref{h:ev:antipodal}. For unit outputs, reflection cancels
\(\mathbb E\langle\widehat S,v\rangle\), whereas
\(\langle\widehat S,u\rangle\le\|u\|\). Cauchy--Schwarz and the
same second-moment identity prove \eqref{h:ev:reflection-expectation}.
\end{proof}

For almost every fixed shared seed, use the prior-preserving Borel
rules justified in Section~\ref{h:stream:section}. Pre-generate all
\(T\) Gaussian rows and grant the learner the full matrix and labels.
After averaging its signal-independent initial state and fresh choices,
its stopped output law is a Borel kernel of those data. Angular error
at most \(\epsilon\) implies Euclidean error at most \(\epsilon\).
If \(T\le d/4\) and \(\epsilon\le1/10\), then
\eqref{h:ev:direct-residual} gives success probability at most
\[
 \frac12+\frac1{8(1-\epsilon^2)}
 \le\frac{62}{99}<\frac23.
\]
This uniform bound also holds for the original conditional success
probabilities and hence for their average over the seed. Thus any
learner with the stipulated average success must satisfy
\begin{equation}\label{h:end:linear-threshold}
 T>d/4.
\end{equation}

\begin{proof}[Proof of Theorem~\ref{h:main}]
Fix the finite constant \(A>0\) in the theorem. Take \(d\) above
the absolute threshold in Proposition~\ref{h:stream:success}, and
large enough that \(a\ge1\) and \(q\ge d/16\). Put
\[
 K_A=\exp(C(1+A))>1,\qquad
 L=\log(1/\epsilon),\qquad B=\lceil T/q\rceil,
\]
where \(C\) is the absolute constant in that proposition. Since
\(M\le Ad^2\) and prior success is at least \(2/3\),
\[
 \frac23\le(2\epsilon K_A^B)^a.
\]
Taking the \(a\)th root gives \(\epsilon K_A^B\ge1/3\).
Since \(B\le T/q+1\),
\begin{equation}\label{h:end:rounding}
 T\ge q\left(\frac{L-\log3}{\log K_A}-1\right).
\end{equation}
Let \(H_A=2(\log3+\log K_A)\). If \(L\ge H_A\), this gives
\[
 T\ge\frac{qL}{2\log K_A}\ge\frac{dL}{32\log K_A}.
\]
If \(L<H_A\), \eqref{h:end:linear-threshold} instead gives
\(T>d/4>dL/(4H_A)\). Thus the same choice
\[
 c_A=\min\left\{\frac1{32\log K_A},\frac1{4H_A}\right\}>0
\]
works uniformly over \(0<\epsilon\le1/10\) and \(M\le Ad^2\)
for this fixed \(A\). No positive constant uniform as \(A\) tends
to infinity is asserted.
\end{proof}

\begin{proof}[Proof of Corollary~\ref{h:subquadratic}]
Since $M(d)=o(d^2)$, there is a dimension $d_M$ after which
$M(d)\le d^2$. Apply Theorem~\ref{h:main} with the numerical choice
$A=1$ at every $d\ge d_M$, so $K_1=e^{2C}$ in the preceding proof.
Its constant $c_1$ is absolute, and the
eventual dimension threshold may depend on the memory sequence
through $d_M$. The argument applies separately at every accuracy
$0<\epsilon(d)\le1/10$, so it requires neither monotonicity of that
sequence nor a relation between memory and accuracy.
\end{proof}

\begin{proof}[Proof of Corollary~\ref{h:worst-signal}]
Average the original fixed-signal experiment over $S\sim\sigma$
before choosing Borel versions. Its measurability makes the resulting
success probability the integral of the fixed-signal success
probabilities. A success probability at least $2/3$ for every
signal therefore gives uniform-prior success at least $2/3$.
Theorem~\ref{h:main} and Corollary~\ref{h:subquadratic} give the
corresponding conclusions with their unchanged constants and
dimension thresholds. Lemma~\ref{h:stream:borel} preserves the
average used in those arguments even if a version changes behavior
at exceptional individual signals.
\end{proof}
